\documentclass[10pt,a4paper]{amsart}
\usepackage[margin=19mm]{geometry}
\usepackage{amsmath,amssymb,mathtools}
\usepackage[scr=rsfs,cal=euler]{mathalfa}
\usepackage{xfrac}
\usepackage[unicode,hidelinks,bookmarks=false]{hyperref}
\newcommand{\ie}{i.e.\ }
\newcommand{\cf}{cf.\ }
\newcommand{\resp}{respectively\ }
\newcommand{\Subsection}{\subsection}
\providecommand{\nobreakdash}{\nobreak\hbox{-}}
\setlength{\emergencystretch}{2em}
\multlinegap0pt
\usepackage{amssymb}
\usepackage{mathtools}
\usepackage[shortlabels]{enumitem}
\usepackage{xcolor}
\usepackage{tikz}
\newtheorem{lemma}{Lemma}
\newcommand{\dist}{\operatorname{dist}}
\title{$p$-Dirichlet spaces over chord-arc domains}
\author{Huaying Wei, Michel Zinsmeister}
\date{Source: arXiv:2410.02183v3 (CC BY 4.0)}
\begin{document}
\maketitle

\noindent\emph{Source and licence.} $p$-Dirichlet spaces over chord-arc domains. Version arXiv:2410.02183v3 (CC BY 4.0), distributed by arXiv under the Creative Commons Attribution 4.0 International licence, http://creativecommons.org/licenses/by/4.0/, as recorded in the arXiv licence metadata for this exact version and shown on the abstract page https://arxiv.org/abs/2410.02183v3. Full licence notice: \texttt{licenses/arxiv-2410.02183-CC-BY-4.0.txt}.\par\smallskip

\noindent\textit{Editorial scope.} Counted unit: the article's lemma 8Kr (source0.tex lines 760--767). The article's complete original proof of it is delivered with it (source0.tex lines 769--819, one \texttt{proof} environment of 51 lines). No mathematics is written, repaired or re-derived: the statement and the proof are the article's own lines, in the article's own notation, and no retained line is substituted or re-flowed.  Retained uncounted as the source-local context the proof needs: lines 708--721; lines 751--754.  Nothing was omitted from any retained span, so the package carries no omission record.  No retained line contains a citation, so the wrapper carries no bibliography and no bibliography entry.  No retained line contains a figure environment, an image-inclusion command or drawing code, so no image is delivered and the package declares no asset dependency.  Editorial formatting only, in addition: an AMS \texttt{amsart} preamble that restates the article's own declarations and macros used by the retained lines, namely the environments \texttt{lemma} on separate counters, together with the standard packages those lines need.  The retained environments print in this document as Lemma 1 in order of appearance, whereas the article prints the same items with the numbering of its own full document.  The complete native source of the article as distributed in the arXiv e-print for this exact version is bundled unchanged as \texttt{sources/166206/source0.tex}.
\begin{lemma}\label{2arc}
Let $\Gamma$ be a rectifiable $K$-quasicircle of length $\ell$.  For
every $x\in\Gamma$ and $R>0$, the set $\Gamma\cap B(x,R)$ is contained in
the union of two 
subarcs $\alpha_+$ and $\alpha_-$ issuing from $x$, such that
\[
  \alpha_+\cup\alpha_-\subset B(x,2KR)
\]
and
\begin{equation}\label{alph}
  {\rm length}(\alpha_+)+{\rm length}(\alpha_-)
  \le2L_x(2KR).
\end{equation}
\end{lemma}

\begin{equation}\label{elem}
  \bigl|\{w:\dist(w,\alpha)<t\}\bigr|
  \le C\bigl(t\ell(\alpha)+t^2\bigr)
\end{equation}

\begin{lemma}\label{8Kr}
Let $1<p<2$ and let $\Gamma$ be a rectifiable $K$-quasicircle.  Then,
for every $x\in\Gamma$ and $r>0$,
\begin{equation}\label{Lx}
  \int_{\Omega\cap B(x,2r)}\delta(w)^{p-2}\,dA(w)
  \leq C_p\bigl(r^{p-1}L_x(8K r)+r^p\bigr). 
\end{equation}
\end{lemma}

\begin{proof}
For $x\in\Gamma$ and $0<t\le2r$, define
\[
  E_t=\{w\in\Omega\cap B(x,2r):\delta(w)<t\}.
\]
If $w\in E_t$, choose $\xi\in\Gamma$ with
$\delta(w)=|\xi-w|$.  Then
\[
  |\xi-x|
  \le|\xi-w|+|w-x|
  =\delta(w)+|w-x|
  <t+2r\le4r.
\]
Lemma~\ref{2arc}, applied with $R=4r$, shows that all such points $\xi$ are
contained in two arcs $\alpha_+$ and $\alpha_-$ lying in
$B(x,8Kr)$ with
\[
  {\rm length}(\alpha_+)+ {\rm length}(\alpha_-)
  \le2L_x(8K r).
\]
Consequently, $E_t$ lies in the $t$-neighborhood of
$\alpha_+\cup\alpha_-$.  By \eqref{elem},
\begin{equation}\label{Et}
  |E_t|\le C\bigl(tL_x(8K r)+t^2\bigr),
  \qquad 0<t\le2r.
\end{equation}

Every $w\in B(x,2r)$ satisfies $\delta(w)\le|w-x|<2r$.  Put
$t_k=2^{1-k}r$ and decompose $\Omega\cap B(x,2r)$ into the layers
\[
  D_k=\{w:t_{k+1}\le\delta(w)<t_k\},
  \qquad k=0,1,2,\cdots.
\]
Since $D_k\subset E_{t_k}$ and, because $p-2<0$,
\[
\delta(w)^{p-2}\le t_{k+1}^{p-2}=2^{2-p}t_k^{p-2}
\quad\text{on }D_k,
\]
estimate \eqref{Et} gives
\begin{align*}
  \int_{\Omega\cap B(x,2r)}\delta(w)^{p-2}\,dA(w)
  &\le C_p\sum_{k=0}^\infty
  t_k^{p-2}\bigl(t_kL_x(8K r)+t_k^2\bigr)\\
  &=C_p\left(
  L_x(8K r)\sum_{k=0}^\infty t_k^{p-1}
  +\sum_{k=0}^\infty t_k^{p}
  \right).
\end{align*}
Both geometric series converge; their sums are bounded by constant
multiples of $r^{p-1}$ and $r^p$, respectively.  This proves \eqref{Lx}.
\end{proof}

\end{document}
